\documentclass[11pt]{article}
\usepackage[a4paper,margin=0.85in]{geometry}
\usepackage{amsmath,amssymb,amsthm,mathtools}
\usepackage{graphicx}
\usepackage{float}
\usepackage[hidelinks]{hyperref}
\usepackage{microtype}

\newtheorem{theorem}{Theorem}
\newtheorem{lemma}[theorem]{Lemma}
\theoremstyle{remark}
\newtheorem{remark}[theorem]{Remark}

\newcommand{\N}{\mathbb N}
\newcommand{\Nzero}{\mathbb N_0}
\newcommand{\vTwo}{v_2}

\title{Explicit weights on $\N$ and $\N^2$ that are not almost monotone}
\author{Candidate full answer to item (4) of arXiv:2407.19982}
\date{26 August 2026}

\begin{document}
\maketitle

\begin{center}
\fbox{\parbox{0.91\textwidth}{\textbf{Status.}
Candidate full answer by explicit counterexample, likely valid, awaiting
specialist review. The source question has an affirmative answer: there are
weights on both $\N$ and $\N^2$ which are not almost monotone.}}
\end{center}

\section{Source question and definitions}

Dabhi calls a map $\omega:\N^2\to[1,\infty)$ a weight when
\[
 \omega(m_1m_2,n_1n_2)
 \leq \omega(m_1,n_1)\omega(m_2,n_2).
\]
It is admissible when
\[
 \lim_{r\to\infty}\omega(m^r,n^r)^{1/r}=1
 \qquad ((m,n)\in\N^2).
\]
It is almost monotone when it is admissible, or when there is $K>0$ such that
\[
 \omega(m_1,n_1)\leq K\omega(m_2,n_2)
\]
whenever $m_1\mid m_2$ or $n_1\mid n_2$ \cite{Dabhi}. For $\N$, the
analogous definitions use $\omega(ab)\leq\omega(a)\omega(b)$,
admissibility along every sequence of powers, and the divisibility relation.

In concluding remark (4) on page 15, the paper asks:
\begin{quote}
``We do not know whether there is a weight $\omega$ on $\N^2$ or $\N$
which is not an almost monotone weight.''
\end{quote}

\begin{figure}[H]
 \centering
 \includegraphics[width=0.98\textwidth]{figures/open_problem_crop.png}
 \caption{The multidimensional definition and item (4) on page 15 of the
 source paper. The earlier $\N^2$ definition on page 2 includes the uniform
 factor $K$ used above.}
\end{figure}

\section{Proof intuition}

The construction separates two requirements that might at first seem to
force each other. A linear term in the exponent of the weight produces
nontrivial spectral growth and hence nonadmissibility. A second, sublinear
term creates unbounded downward jumps along divisibility.

For the second term, write an integer in base 4 and evaluate the same digit
string in base 2. This digit cost is subadditive: when two base-4 expansions
are added, every carry replaces four copies of cost $2^i$ by one copy of cost
$2^{i+1}$, so the cost decreases. Nevertheless, just before a power of 4 all
digits are 3, while at the power of 4 only one digit is nonzero. This creates
an arbitrarily large downward jump. Combining the linear term with this digit
cost gives exactly the two properties needed.

\section{The digit-cost lemma}

For $r\in\Nzero$, let
\[
 r=\sum_{i=0}^J d_i4^i,\qquad d_i\in\{0,1,2,3\},
\]
be its base-4 expansion, and define
\[
 c(r):=\sum_{i=0}^Jd_i2^i.
\]

\begin{lemma}\label{lem:cost}
The function $c:\Nzero\to\Nzero$ is subadditive. Moreover,
\[
 \frac{c(r)}r\longrightarrow0,
 \qquad
 c(4^j-1)=3(2^j-1),
 \qquad
 c(4^j)=2^j.
\]
\end{lemma}

\begin{proof}
Place the base-4 digits of two integers on top of each other. Before carrying,
the digit cost of their sum is exactly the sum of their two digit costs. A
carry at position $i$ replaces four units at that position by one unit at
position $i+1$. Its contribution to the digit cost therefore changes from
$4\cdot2^i$ to $2^{i+1}$ and cannot increase. After all carries, the resulting
cost is $c(r+s)$, proving
\[
 c(r+s)\leq c(r)+c(s).
\]

If $4^J\leq r<4^{J+1}$, then
\[
 c(r)\leq3\sum_{i=0}^J2^i=3(2^{J+1}-1),
\]
whereas $r\geq4^J$. Thus $c(r)/r\leq6/2^J$, which tends to zero.
Finally, $4^j-1$ has $j$ base-4 digits, all equal to 3, while $4^j$ has a
single digit 1 at position $j$. The two exact formulas follow.
\end{proof}

\section{Full answer}

Let $\vTwo(n)$ denote the exponent of 2 in the prime factorization of $n$,
with $\vTwo(1)=0$.

\begin{theorem}\label{thm:answer}
Define
\[
 \phi(r):=r+c(r)\quad(r\in\Nzero),
 \qquad
 \boxed{\ \omega(n):=2^{\phi(\vTwo(n))}\quad(n\in\N).\ }
\]
Then $\omega$ is a weight on the multiplicative semigroup $\N$, but it is not
almost monotone. Furthermore,
\[
 \Omega(m,n):=\omega(m)
\]
is a weight on $\N^2$ which is not almost monotone. Consequently item (4) of
the source paper has an affirmative answer for both $\N$ and $\N^2$.
\end{theorem}

\begin{proof}
The function $\phi$ is nonnegative and subadditive by
Lemma~\ref{lem:cost}. Since $\vTwo(ab)=\vTwo(a)+\vTwo(b)$,
\[
 \omega(ab)
 =2^{\phi(\vTwo(a)+\vTwo(b))}
 \leq2^{\phi(\vTwo(a))+\phi(\vTwo(b))}
 =\omega(a)\omega(b).
\]
Thus $\omega:\N\to[1,\infty)$ is a weight.

It is not admissible. Indeed, Lemma~\ref{lem:cost} gives
\[
 \lim_{r\to\infty}\omega(2^r)^{1/r}
 =\lim_{r\to\infty}2^{(r+c(r))/r}=2>1.
\]
It remains to rule out the other branch of almost monotonicity. For $j\geq1$,
put
\[
 a_j=2^{4^j-1},\qquad b_j=2^{4^j}.
\]
Then $a_j\mid b_j$, while the exact formulas in Lemma~\ref{lem:cost} give
\begin{align*}
 \phi(4^j-1)-\phi(4^j)
 &= (4^j-1)+3(2^j-1)-(4^j+2^j)\\
 &=2^{j+1}-4.
\end{align*}
Hence
\[
 \frac{\omega(a_j)}{\omega(b_j)}
 =2^{2^{j+1}-4}\longrightarrow\infty.
\]
No constant $K$ can satisfy $\omega(a)\leq K\omega(b)$ whenever $a\mid b$.
Because $\omega$ is also nonadmissible, it is not almost monotone.

For $\Omega(m,n)=\omega(m)$, submultiplicativity follows immediately from
that of $\omega$. It is nonadmissible because
\[
 \lim_{r\to\infty}\Omega(2^r,1)^{1/r}=2.
\]
Finally, $(a_j,1)$ divides $(b_j,1)$ coordinatewise and
\[
 \frac{\Omega(a_j,1)}{\Omega(b_j,1)}
 =\frac{\omega(a_j)}{\omega(b_j)}\longrightarrow\infty.
\]
Thus $\Omega$ fails the uniform almost-monotonicity inequality even on this
one-coordinate chain. This proves the claim on $\N^2$.
\end{proof}

\section{Verification, scope, and novelty check}

The proof has three independent audit points. First, base-4 carrying proves
subadditivity of $c$. Second, $c(r)=o(r)$ makes the linear term determine the
nonadmissible spectral growth. Third, the explicit pair of consecutive
exponents $4^j-1,4^j$ gives an unbounded downward ratio along divisibility.
The lift to $\N^2$ preserves all three properties.

A bounded exhaustive check verified $c(r+s)\leq c(r)+c(s)$ for all
$0\leq r,s\leq1000$ and verified the exact jump formula for $1\leq j\leq12$.
This is only a sanity check; the carry argument is the proof.

The page-15 multidimensional restatement mentions $K>0$ but omits $K$ from
the following displayed comparison, whereas the formal $\N^2$ definition on
page 2 includes it. The construction above has ratios tending to infinity,
so it fails both the literal monotone comparison and the more permissive
uniform-$K$ comparison. The disjunction ``admissible or uniformly monotone''
is also handled explicitly: the weight fails both branches.

A bounded novelty check through 26 August 2026 searched the run's registry,
solution, attempt, and proof-gap indexes; the exact phrase ``not an almost
monotone weight''; the source title, author, and arXiv id; and arXiv searches
combining ``almost monotone'', ``weight'', and ``Dirichlet''. The search found
the source preprint and its 2026 journal publication, but no independent paper
answering item (4). The construction is elementary, so this search does not
establish historical priority beyond those bounds. Novelty confidence is
therefore moderate, while mathematical confidence is high.

The highest-value specialist check is to verify the carry argument and the
source-definition match. Under either natural reading of the page-15 wording,
the proof is complete.

\begingroup
\small
\begin{thebibliography}{9}
\bibitem{Dabhi}
P. A. Dabhi,
\emph{On inversion of absolutely convergent weighted Dirichlet series in two
variables}, arXiv:2407.19982 (2024); Mathematica Bohemica, first online
27 January 2026.
\end{thebibliography}
\endgroup

\end{document}
